\subsection{Operaciones no conmutativas}

\begin{frame}
  \frametitle{Consultas no conmutativas}

  El algoritmo de consultas en caminos funciona también para operaciones \textbf{no conmutativas} (por ejemplo, producto de matrices).

  \vspace{0.75em}

  \begin{itemize}
    \item Guardamos pares con la suma de izquierda a derecha y la suma de derecha a izquierda.
    \item En cada paso del algoritmo elegimos la dirección correcta según cuál punta del camino estamos resolviendo.
  \end{itemize}
\end{frame}

\begin{frame}[fragile]
  \frametitle{Consultas no conmutativas: estructuras}

\begin{lstlisting}
struct Mono {
    static Mono neut() { /* ... */ }
};
Mono operator + (Mono x, Mono y);

struct Conm {
    Mono ltr, rtl;
    static Conm make(Mono x) { return {x, x}; }
    static Conm neut() { return {Mono::neut(), Mono::neut()}; }
};
Conm operator + (Conm x, Conm y) {
    return {x.ltr + y.ltr, y.rtl + x.rtl};
}
\end{lstlisting}
\end{frame}

\begin{frame}[fragile]
  \frametitle{Consultas no conmutativas: \texttt{query}}

\begin{lstlisting}
Mono query(int u, int v) {
  int tu = htop[u], tv = htop[v];
  if (tu == tv) {
    if (dep[u] < dep[v]) {
      return tree[tu].query(dep[u]-dep[tu],dep[v]-dep[tu]+1).ltr;
    } else {
      return tree[tv].query(dep[v]-dep[tv],dep[u]-dep[tv]+1).rtl;
    }
  } else {
    if (dep[tu] > dep[tv]) {
      return tree[tu].query(0, dep[u]-dep[tu]+1).rtl
           + query(p[tu], v);
    } else {
      return query(u, p[tv])
           + tree[tv].query(0, dep[v]-dep[tv]+1).ltr;
    }
  }
}
\end{lstlisting}
\end{frame}
